\title{Parameter-Efficient Orthogonal Finetuning via Factorization}
\section{Concluding Remarks and Limitations}

Originally, the butterfly structure was developed as part of the Cooley-Tukey algorithm to facilitate the fast Fourier transform. In the context of the Fourier transform, modifications to a single frequency can propagate and induce changes across the entire spatial domain, mirroring our scenario for information transmission in which each node from an initial layer can send information to all nodes in the final layer. The butterfly structure has also become an established design for efficient data exchange in computer network topology~\cite{solihin2015fundamentals,li2011linear}. Considering $k \geq 2$ as a power of $2$, we first introduce the \emph{butterfly factor} $\bm{B}^F(k)$, defined as

To investigate this problem, we introduce the approach of expressing a dense orthogonal matrix as a sequence of products involving several sparse orthogonal matrices. To establish a systematic and accessible way of analyzing sparsity structures in orthogonal matrix decompositions, we recast the task of constructing a dense orthogonal matrix in OFT as an instance of information transmission. Concretely, forming a dense matrix $\bm{R}\in\mathbb{R}^{d\times d}$ through the multiplication of $m$ square matrices, $\thickmuskip=2mu \medmuskip=2mu \bm{R}=\bm{B}_m\bm{B}_{m-1}\cdots\bm{B}_1$, can be conceptualized as propagating information through a $d\times (m+1)$ node lattice, depicted in Figure~\ref{fig:info_trans}. This analogy stems from the fact that a dense $d$-dimensional square matrix can represent fully connected paths from one group of $d$ nodes to another. In this framework, a zero entry $R_{ij}$ in $\bm{R}$ implies that transmission from the $j$-th to the $i$-th node is blocked, whereas a nonzero $R_{ij}$ signifies an active path. Consequently, expressing $\bm{R}$ as a product $\bm{B}_m\bm{B}_{m-1}\cdots\bm{B}_1$ corresponds to a sequence of information exchanges, dictated by the connection patterns in each $\bm{B}_i$. The transmission process follows the sequence, initiating with $\bm{B}_1$ and concluding with $\bm{B}_n$. For illustration, consider Figure~\ref{fig:info_trans}, which examines the decomposition $\bm{R}=\bm{B}_5\bm{B}_4\bm{B}_3\bm{B}_2\bm{B}_1$, showing their sparsity patterns and the corresponding transmission graph. Here, the displayed graph unpacks the multiplication across the factors, interpreting each $\bm{B}_i$ as a map from nodes at level $i$ to those at level $(i+1)$. Precisely, the $(j_1,j_2)$ entry in $\bm{B}_i$ reflects the existence of a directed edge from the $j_2$-th node at level $i$ to the $j_1$-th node at level $(i+1)$, with a zero denoting no direct link. To achieve a dense matrix via $\bm{B}_5\bm{B}_4\bm{B}_3\bm{B}_2\bm{B}_1$, every node in the initial layer must be able to transmit information to all nodes in the final, $6$-th, layer. However, examining $\bm{R}=\bm{B}_4\bm{B}_3\bm{B}_2\bm{B}_1$, which concerns information paths from the first to the fifth level, reveals that node $1$ is unable to reach node $3$, implying that $(3,1)$ in $\bm{B}_4\bm{B}_3\bm{B}_2\bm{B}_1$ must be zero. Similar associations between graph connectivity and the sparsity structure are observed for $\bm{R}=\bm{B}_3\bm{B}_2\bm{B}_1$ and $\bm{R}=\bm{B}_2\bm{B}_1$. In summary, for a matrix $\bm{R}\in\mathbb{R}^{d\times d}$ to possess full density, the factors $\bm{B}_m,\cdots,\bm{B}_1$ must establish directed edges in a $d\times (m+1)$ grid, where each edge only links nodes in adjacent levels (i.e., columns), and ensure that all nodes at the initial level have paths reaching all nodes at the $(m+1)$-th level.

Formally, OFT introduces a trainable orthogonal matrix $\bm{R}\in\mathbb{R}^{d\times d}$ applied to a pretrained linear layer $\bm{W}^0\in\mathbb{R}^{d\times n}$. This changes the forward computation from $\bm{z}=(\bm{W}^0)^\top\bm{x}$ to $\bm{z}=(\bm{R}\bm{W}^0)^\top\bm{x}$, where $\bm{x}\in\mathbb{R}^{d}$ and $\bm{z}\in\mathbb{R}^{n}$ represent input and output vectors, respectively. OFT ensures that $\bm{R}$ is initially the identity matrix, guaranteeing zero initialization. To maintain orthogonality during finetuning, we adopt the Cayley parameterization following \cite{qiu2023controlling,liu2021orthogonal}: $\bm{R}=(\bm{I}+\bm{Q})(\bm{I}-\bm{Q})^{-1}$, with $\bm{Q}$ being a skew-symmetric matrix so that $\bm{Q}=-\bm{Q}^\top$. For parameter efficiency, $\bm{R}$ is parameterized as a block-diagonal matrix: $\text{diag}(\bm{R}_1,\bm{R}_2,\cdots,\bm{R}_r)$, where each block $\bm{R}_i\in\mathbb{R}^{b\times b}$ is an orthogonal matrix and $br = d$. However, this parameter-efficient structure requires splitting the dimensions of a neuron (i.e., the column vectors in $\bm{W}^0$) into $r$ separate groups, with each group independently transformed by a different orthogonal matrix. Despite the demonstrated empirical effectiveness of such block-diagonal constructions, it is generally arbitrary to segment the neuron’s dimensions solely by index, suggesting that a dense orthogonal matrix could be preferable. This naturally leads to the following question: \emph{Is it possible to design a dense orthogonal matrix while preserving parameter efficiency?}

The central challenge lies in constructing a dense orthogonal matrix while maintaining a low parameter count. Traditionally, representing a $d$-dimensional dense orthogonal matrix demands $\mathcal{O}(d^2)$ parameters, making this goal appear unattainable. However, we introduce an innovative strategy: composing the dense orthogonal matrix as a product of several factorized sparse matrices. This idea leverages the realization that parameter savings are possible if additional computational overhead per iteration is acceptable. Since the orthogonal matrix is formed through multiplication of sparse factors, repeated multiplications are required during each round of training. To provide an alternative viewpoint, we model the creation of a dense orthogonal matrix as a problem of information flow. In this analogy, generating such a matrix via products of sparse factors is akin to disseminating information across a grid-like network. Framing the task in this way uncovers diverse approaches to sparse matrix factorization, enabling significant reductions in the parameter count while preserving the flexibility needed to span dense orthogonal matrices.

\begin{equation}\label{eq:but_factor}
\begin{aligned}
    \bm{B}^F(k)=\begin{bmatrix}
    \text{diag}(\bm{d}_1) & \text{diag}(\bm{d}_2) \\
    \text{diag}(\bm{d}_3) & \text{diag}(\bm{d}_4)
    \end{bmatrix}\in\mathbb{R}^{k\times k},
\end{aligned}
\end{equation}

\section{Intriguing Insights and Discussions}

\section{Orthogonal Parameterization by Butterfly Factorization}

Here, we adopt the shorthand $\tilde{\bm{B}}^b_{(d,i)}$ to refer to $\tilde{\bm{B}}^b(d,2^{m - i+1})$ for ease of notation, with $\bm{I}_d$ denoting the $d \times d$ identity matrix. The orthogonal matrix $\bm{R}(m,b)\in\mathbb{R}^{d\times d}$ is constructed by sequentially multiplying $m$ orthogonal butterfly blocks. For notational clarity, we refer to the variant using $\bm{R}(m,\frac{b}{2})$ as BOFT($m$,$b$), requiring $b\geq 2$. Setting $m=1$ yields BOFT($1$,$b$), which is equivalent to the block-diagonal OFT~\cite{qiu2023controlling} parameterized by block size $b$. Taking $b=d$ in BOFT($1$,$d$) recovers the unconstrained full orthogonal matrix of the standard OFT~\cite{qiu2023controlling}. When choosing $m=\log \frac{2d}{b}$ and $b$ fixed, the block butterfly matrix $\bm{B}^{\frac{b}{2}}(d)$ becomes $\bm{R}$, producing a dense orthogonal transform. More broadly, BOFT($m$,$b$) utilizes $\frac{1}{2}(b-1)dm$ total parameters when adapting a linear transformation with dimensions $d\times n$. In the special case where a butterfly structure is employed, i.e., $m=\log d$ and $b=2$, the required parameter count is $\mathcal{O}(d\log d)$. This contrasts with the original OFT (full orthogonal matrix), which scales as $\mathcal{O}(d^2)$, and the block-diagonal OFT (with $r$ blocks) requiring $\mathcal{O}(bd)$. Hence, to form a dense orthogonal transform, the original OFT must use $b=d$, whereas BOFT is flexible in the choice of $b$.

\begin{theorem}[Expressivity of BOFT]\label{thm:exp_boft}
    BOFT possesses greater expressivity than OFT when the block size is held constant. To enable the butterfly matrix to approximate all orthogonal matrices of dimension $d$, one can take products of butterfly matrices in the form $\bm{B}_{d-1,1}(d)\bm{B}^\top_{d-1,2}(d)\cdots\bm{B}_{1,1}(d)\bm{B}^\top_{1,2}(d)$, where each $\bm{B}_{i,j}(d),\forall i,\forall j$ is a butterfly matrix.
\end{theorem}

\textbf{Identity Initialization Strategy for BOFT.} In transfer learning settings, it is standard practice to initialize adaptation layers such that the model's predictions remain close to those of the pretrained model; LoRA achieves this by setting low-rank adaptation weights to zero. Consistently, BOFT initializes all butterfly components to identity (which, via the Cayley transform, amounts to zero-initializing the underlying skew-symmetric matrices).

We begin by outlining foundational aspects of OFT. In order to finetune the pretrained weight matrix, OFT redefines the updated weight matrix as the multiplication of a trainable orthogonal matrix with the fixed, pretrained matrix. Unlike LoRA—which modifies the weights through the addition of a low-rank matrix—OFT employs a multiplicative orthogonal transformation to update the weights. While LoRA attains parameter efficiency via a low-rank configuration, the initial design of OFT achieves this using a block-diagonal orthogonal matrix structure~\cite{qiu2023controlling}; specifically, reducing the block size further enhances OFT’s parameter efficiency. Figure~\ref{fig:comp_lora} offers a clear visual comparison. The central reasoning behind employing orthogonal transformations for finetuning lies in the ability to maintain neuron pair-wise angles~\cite{liu2018learning,liu2021orthogonal,qiu2023controlling}, which in turn supports the retention of semantic information from the pretraining phase. 

\textbf{Inductive bias and generalization in BOFT}. In BOFT, the transformation $\bm{R}(m,b)$ often restricts the solution to a specific structured subset of the orthogonal group, thus narrowing the hypothesis space and instilling an inherent inductive bias. It is posited that the particular inductive bias conferred by butterfly factorization is advantageous for model generalization, given its underlying regularities shared with widely-used linear transforms such as the discrete Fourier, sine/cosine, and Hadamard transforms~\cite{dao2019kaleidoscope}. Additionally, the structured sparsity in BOFT’s matrix factorization may yield further implicit regularization effects~\cite{gunasekar2017implicit,li2018algorithmic,liu2019neural}.

Model finetuning fundamentally aims to maintain similarity between the pretrained and the adapted models according to specific metrics, thereby retaining the pretrained model’s acquired knowledge. A common practice in modern finetuning is to employ a low learning rate, an intuitive strategy that helps to keep the changes between the pretrained and finetuned models minimal. Beyond this heuristic, the inherent structure of weight matrices has motivated more principled strategies that seek to maintain the relationships among neurons, particularly by conserving the pairwise angular relationships in the weight matrices. This concept underpins Orthogonal Finetuning~(OFT)~\cite{qiu2023controlling}, a novel finetuning methodology where each neuron within a layer undergoes transformation by a shared orthogonal matrix, guaranteeing preservation of pairwise neuron angles during adaptation. Although OFT has been shown to deliver notable performance and stability in finetuning text-to-image diffusion models~\cite{qiu2023controlling}, a major drawback is the substantial number of parameters involved, stemming from the large size of orthogonal matrices. To overcome this limitation, OFT proposes utilizing a block-diagonal configuration, which significantly reduces parameter count. Nevertheless, this gain in parameter efficiency comes at a cost: the block-diagonal orthogonal matrix imposes a predetermined sparsity pattern and segments the transformation process, thereby reducing expressivity and limiting its capacity to represent classical linear operations. Additionally, the challenge of determining an optimal sparsity pattern for the block-diagonal structure remains unresolved.

Interpreting OFT through the lens of information transmission offers insights into why the factorization matrices display specific sparsity patterns. In Figure~\ref{fig:blk_diag}, the original OFT’s matrix $\bm{R}$ is depicted with a block-diagonal form, which helps decrease the number of parameters to be optimized. However, this architecture inherently limits the capacity to realize a dense $\bm{R}$. Our objective is to synthesize a dense orthogonal matrix by composing $m$ sparse orthogonal factors, aiming to use as few effective trainable parameters as necessary. The information transmission perspective translates this objective into two main requirements: \emph{(i)} \textbf{dense connectivity}—each initial-level node should be able to reach every final-level node via at least one path; and \emph{(ii)} \textbf{minimal free edges}—subject to orthogonality, the edge count must be kept to the lowest feasible value. Importantly, enforcing orthogonality imposes strict limitations on edges connecting adjacent layers. Specifically, each $\bm{B}_i$ must be full-rank—which is essential for orthogonality—so a bijective mapping involving $d$ edges must connect the nodes of consecutive layers, resulting in no fewer than $d$ edges (for example, Figure~\ref{fig:info_trans} illustrates this with $d=4$). Such necessary edges are essential for preserving orthogonality and do not represent trainable degrees of freedom, as they involve fixed (non-trainable) values such as $\pm 1$ in a $d\times d$ orthogonal matrix with $d$ non-zero entries. Because the orthogonality condition introduces constraints that tie parameters together more tightly than in unconstrained full matrices, certain structural dependencies arise. For instance, if a $2\times 2$ orthogonal matrix has a zero at position $(1,1)$, this may enforce another zero at $(2,2)$, so omitting one connection may necessitate omitting another. Thus, the allowable arrangement of edge connections is shaped by the orthogonality requirement, sometimes compelling additional removals or additions. Within this context, focusing on the pattern of non-zero entries in the resulting orthogonal matrix proves especially helpful: the information transmission approach aids analysis in OFT because it emphasizes the importance of which entries of $\bm{R}$ are trainable (non-zero), while the actual values of these entries are of secondary concern.

In this work, we introduce Orthogonal Butterfly, a versatile, parameter-efficient finetuning technique for foundation models that leverages the butterfly architecture. The core principle underlying our approach is that increased parameter-efficiency can be attained by representing a dense orthogonal matrix as a product of several sparse orthogonal matrices. To facilitate the identification of suitable matrix decompositions, we present a graphical framework grounded in information transmission principles. Within this paradigm, we demonstrate that the butterfly structure provides a practical and effective solution for the sparse factorization of orthogonal matrices. Our empirical studies showcase the robustness of BOFT in finetuning a wide array of models, including large language models, vision models, and text-to-image generative models, affirming its effectiveness as a general-purpose finetuning technique.

\textbf{Orthogonal finetuning as learning bilinear similarity}. BOFT can be reformulated as optimizing a bilinear similarity expression of the form $\bm{w}^0_i\bm{R}\bm{x}$, where $\bm{w}^0_i$ is the $i$-th column of $\bm{W}^0$, interpreted as an individual neuron. In this perspective, BOFT’s task is to adapt the bilinear similarity matrix $\bm{R}$, subject to a stringent orthogonality constraint. This formulation directly relates to problems in distance metric learning~\cite{xing2002distance} and aligns with the theory of bilinear forms~\cite{roman2005advanced}.

Nevertheless, BOFT has its drawbacks. The resulting orthogonal matrix in BOFT emerges from the multiplication of multiple orthogonal matrices, which introduces a marginally higher training overhead relative to OFT. The challenge of reducing BOFT's training computational cost remains unsolved. On the positive side, after finetuning, the orthogonal matrices obtained by BOFT can be merged directly into the pretrained model parameters without imposing any extra inference-time delay. Additionally, it is not yet established whether the butterfly network constitutes the optimal architecture for information propagation. Our proposed information transmission framework opens avenues for insights from other disciplines—particularly computer networking, where topological efficiency for information transfer is a major focus. This cross-disciplinary perspective suggests that even more effective network designs might be developed for constructing dense orthogonal matrices.

\section{An Information Transmission View on Orthogonal Finetuning}

\textbf{Comparison to butterfly-based sparse training}. Several prior studies~\cite{dao2019learning,dao2019kaleidoscope,chen2022pixelated,dao2022monarch} have explored sparse training approaches using butterfly parameterization. These works typically involve directly expressing the weight matrices using butterfly structures and training networks from the beginning. In contrast, \cite{dao2022monarch} introduces a method for finetuning existing models by initially projecting pretrained weights onto a modified butterfly matrix, followed by optimization of these projected components for specific tasks. BOFT adopts a fundamentally distinct approach to finetuning, where transformations are applied to the weights using weight matrices that are shared across layers.

\begin{itemize}[leftmargin=*,nosep]
    \item We analyze the challenge of achieving parameter-efficient orthogonal finetuning by introducing a novel information transmission framework, which reframes the task of designing a parameter-efficient dense orthogonal matrix as a communication process within a graph-based structure.
    \item Building on insights from butterfly networks in the Cooley-Tukey algorithm, we introduce Orthogonal Butterfly, a new method for orthogonal finetuning with enhanced parameter efficiency, formulated within our information transmission perspective.
    \item We offer several theoretical explanations elucidating how BOFT manages to considerably decrease the number of trainable parameters while preserving expressive capacity and ensuring stable training dynamics. Additionally, the matrix factorization inherent to BOFT naturally supports an interesting mechanism for weight interpolation (refer to Figure~\ref{fig:boft_interpolation}).
    \item We are the first to extend orthogonal finetuning to a broader range of applications beyond controllable text-to-image generation~\cite{qiu2023controlling}, demonstrating its versatility as a general-purpose finetuning tool. In our experiments, we deploy BOFT across diverse adaptation tasks, including those in computer vision and natural language processing, where it achieves significant gains over leading alternatives, thus highlighting its advanced parameter efficiency and robust generalization properties.
\end{itemize}

\textbf{Expressivity of BOFT}. The butterfly architecture, when combined with permutations, is capable of exactly reconstructing various well-known fast linear transforms~\cite{dao2019learning,dao2019kaleidoscope} (\emph{e.g.}, the fast Fourier transform and the Hadamard transform). Nevertheless, the extent to which our orthogonal butterfly matrix can approximate arbitrary orthogonal matrices has not yet been fully investigated. To examine this, we perform a simulation in which a randomly sampled dense orthogonal matrix~\cite{anderson1987generation} of dimension $512\times 512$ is approximated. In Figure~\ref{fig:para_eff}, we present results averaged over 10 random seeds. The vertical axis reports the approximation error, while the horizontal axis corresponds to the number of effective trainable parameters. Curves of the same color represent BOFT models with identical block sizes, and the leftmost point for each corresponds to BOFT($1$,$b$) (that is, the original block-diagonal OFT with block size $b$). Generally, BOFT demonstrates superior parameter efficiency relative to OFT. For instance, BOFT($9$,$2$) offers greater expressivity than BOFT($1$,$16$) while utilizing far fewer parameters. Using smaller $b$ and larger $m$ in BOFT tends to yield greater parameter efficiency; for example, BOFT($6$,$4$) achieves a comparable approximation error to BOFT($2$,$16$) with significantly fewer parameters. Overall, the butterfly matrix represents a more constrained and structured subset within the orthogonal group compared to a block-diagonal design, resulting in BOFT being theoretically more expressive than OFT for a fixed block size.

\begin{equation}
\begin{aligned}
    \!\!\bm{B}(d)=\tilde{\bm{B}}(d,d)\!\cdot\!\begin{bmatrix}
    \bm{B}_1(\frac{d}{2})\!\!\!\!\! & \bm{0} \\
    \bm{0}\!\!\!\!\! &\bm{B}_2(\frac{d}{2})
    \end{bmatrix}= \tilde{\bm{B}}(d,d)\tilde{\bm{B}}(d,\frac{d}{2})\cdots\tilde{\bm{B}}(d,2),
\end{aligned}
\end{equation}

Recent advancements exemplified by models such as ChatGPT~\cite{brown2020language,chen2021evaluating} and Stable Diffusion~\cite{rombach2022high} have highlighted the outstanding capacity for generalization within large foundation models. Notably, their performance gains are closely tied to dramatic increases in parameter count

Within this information transmission paradigm, we take inspiration from the butterfly networks found in the Cooley-Tukey fast Fourier transform algorithm~\cite{cooley1965algorithm}, which allow efficient exchange of information between network nodes~\cite{klasing1994broadcasting}. In particular, the butterfly topology inherent to Cooley-Tukey suggests a specific sparse matrix decomposition—termed butterfly factorization. Employing butterfly factorization allows for the generation of a dense $d$-dimensional matrix through a sequence of $\mathcal{O}(\log d)$ sparse matrices, where each contains only $\mathcal{O}(d)$ nonzero entries. By ensuring each such factor is orthogonal, we obtain an efficient orthogonal parameterization that requires merely $\mathcal{O}(d\log d)$ parameters—an exponential improvement over the full dense case. Building on this foundation, we present Orthogonal Butterfly~(BOFT), a new approach for parameter-efficient finetuning. BOFT generalizes OFT as a special instance, establishing a flexible and broadly applicable orthogonal finetuning scheme. Both block-diagonal and butterfly-based approaches share a unifying theme: exploiting the sparsity structure to restrict the number of learnable parameters in orthogonal matrices. A comparison to the low-rank parameterizations in LoRA~\cite{hulora2022} is also insightful; both low-rank and sparse matrices are principal examples of structured matrices~\cite{candes2011robust}, offering efficient solutions for reducing the parameter budget.

\textbf{Parameter-efficient finetuning (PEFT)}. With the advent of increasingly large and capable foundation models (\emph{e.g.}, \cite{brown2020language,radford2021learning,rombach2022high,kirillov2023segment}), efficiently adapting these models to specific downstream tasks has emerged as a topic of significant research interest~\cite{treviso2023efficient,ding2023parameter,lialin2023scaling}. Among a diverse array of PEFT strategies~\cite{houlsby2019parameter,aghajanyan2020intrinsic,hulora2022,edalati2022krona,wang2022adamix,gheini2021cross,zaken2022bitfit,guo2020parameter,sung2021training,ansell2022composable,lester2021power,li2021prefix,vu2022spot,he2021towards,mao2021unipelt,karimi2021compacter,liu2022few,sung2022lst,chen2022parameter,jia2022visual,chen2022adaptformer,zhang2022neural,jie2023fact,lian2022scaling,luo2023towards}, reparameterization-based approaches~\cite{aghajanyan2020intrinsic,hulora2022,edalati2022krona} are particularly pertinent to the context of this study. LoRA~\cite{hulora2022}, for instance, introduces additional low-rank matrices to the pretrained weights, updating the weight matrix through the addition of a low-rank product, and achieves strong empirical results for natural language applications. Recognizing that LoRA uses a fixed rank for its added matrices across all layers, subsequent methods~\cite{zhang2022adaptive,valipour2022dylora,zhang2023increlora} have explored dynamically varying the rank per layer to better allocate the parameter budget. The straightforward nature of low-rank weight reparameterization has thus resulted in widespread adoption~\cite{dettmers2023qlora,zi2023delta,chavan2023one}. In parallel, \cite{qiu2023controlling}, building on the insight that hyperspherical energy relates to generalization~\cite{liu2018learning,liu2021orthogonal}, put forward orthogonal finetuning (OFT) as a compelling alternative for finetuning text-to-image diffusion models. OFT introduces an orthogonal transformation within each layer, facilitating better generalization and more reliable training outcomes when compared to LoRA, albeit typically with a higher number of tunable parameters. Consequently, further increasing the parameter efficiency of OFT remains an important objective. Additionally, the broader applicability of OFT to adaptation tasks outside of text-to-image diffusion~\cite{qiu2023controlling} has yet to be thoroughly examined. To this end, BOFT employs butterfly factorization to enhance the parameter efficiency of OFT. As a result, it becomes possible to empirically verify the advantages of orthogonal finetuning for a broader range of adaptation scenarios.

where each $\bm{d}_i \in \mathbb{R}^{\frac{k}{2}}$ for $i=1,\ldots,4$ are vector parameters. For $d=2^N$, we then define a $d$-dimensional \emph{butterfly matrix} $\bm{B}(d)\in\mathbb{R}^{d\times d}$ recursively as follows:


\begin{abstract}
While large foundation models have become increasingly prevalent, the substantial costs associated with training them from the ground up make this approach impractical. Consequently, the development of methods to efficiently tailor these advanced models to specific downstream tasks has become crucial. This work focuses on a structured finetuning approach known as Orthogonal Finetuning (OFT), which is designed for adapting models to new tasks. Although OFT is effective in terms of generalization, it tends to rely on a significant number of trainable parameters due to the inherent dimensionality of orthogonal matrices. To overcome this issue, we reinterpret OFT through the lens of information transmission, extracting several essential criteria that could enhance parameter efficiency. Taking inspiration from the efficient communication enabled by the Cooley-Tukey fast Fourier transform, we introduce a more compact orthogonal parameterization based on butterfly matrices. Incorporating this parameterization into OFT yields a new, parameter-efficient finetuning technique termed Orthogonal Butterfly (BOFT). BOFT generalizes OFT—of which it is a strict extension—and offers a broader framework for orthogonal finetuning. We thoroughly evaluate this method by adapting large vision transformers, language models, and text-to-image diffusion architectures across multiple downstream vision and language benchmarks.
\end{abstract}

\input{rebuttal-results}

The implication of Theorem~\ref{thm:exp_boft} is that BOFT can be naturally generalized—specifically, the resultant orthogonal matrix can be parameterized as $\thickmuskip=2mu \medmuskip=2mu \bm{R}^G(m_1,b_1,m_2,b_2,l)=\bm{R}_{l,1}(m_1,b_1)\bm{R}_{l,2}^T(m_2,b_2)\cdots\bm{R}_{1,1}(m_1,b_1)\bm{R}_{1,2}^T(m_2,b_2)$, with $\bm{R}_{i,j}^T(m,b)$ referring to the orthogonal matrices utilized within BOFT. For the special case where $m_1=m_2=\log d$, $b_1=b_2=2$, and $l=d-1$, the construction $\bm{R}^G(m_1,b_1,m_2,b_2,l)$ is sufficient to capture the entire orthogonal group. This type of composition is referred to as the kaleidoscope hierarchy~\cite{dao2019kaleidoscope}. Nonetheless, it is important to recognize that higher expressivity does not guarantee better finetuning outcomes—unrestricted finetuning, despite offering full expressivity, can often lead to poor performance. Achieving a suitable balance between expressivity and imposed structure (regularity) is central to ensuring model finetuning generalizes well. BOFT expands the space of finetuning parameters through its structured priors, providing an improved balance between expressivity and regularization.

\end{document}

\textbf{Spectral properties}. Orthogonal finetuning typically achieves superior preservation of spectral properties compared to LoRA, owing to its exact maintenance of the spectral norm of the original pretrained weight matrix $\bm{W}^0$. This invariance can be understood by examining the singular value decomposition $\bm{W}^0=\bm{U}\bm{\Sigma}\bm{V}^\top$, where $\bm{U}$ and $\bm{V}$ are orthogonal matrices and $\bm{\Sigma}$ contains the singular values on its diagonal. The approaches OFT and BOFT both left-multiply $\bm{W}^0$ by an orthogonal matrix $\bm{R}$, leading to finetuned parameters $\bm{R}\bm{U}\bm{\Sigma}\bm{V}^\top$. This transformation leaves the dominant singular value, i.e., the spectral norm, unchanged. Retaining this property has been empirically linked to improved training robustness and enhanced generalization capabilities~\cite{miyato2018spectral,yoshida2017spectral}. Further theoretical discussions and mathematical results can be found in Appendix~\ref{app:sn_property}.

A straightforward implementation of a fully connected architecture between two layers requires $\mathcal{O}(d^2)$ connections (\emph{i.e.}, represented by a single dense orthogonal matrix), which corresponds to $d^2 - d$ edges (for $d=4$, this totals 12 edges). As illustrated in Figure~\ref{fig:info_trans}, a particular example of matrix factorization involves just $10$ edges overall, which is fewer than those needed for a single dense orthogonal matrix. This conceptual framework offers a graphical lens through which to analyze the parameter efficiency of OFT, and provides an accessible means to construct viable factorizations. We are motivated by an intriguing topology from the Cooley-Tukey algorithm, specifically the butterfly graph~\cite{cooley1965algorithm}, which achieves dense connectivity between $d$ input and $d$ output nodes using only $\mathcal{O}(d\log d)$ edges. As an example, the network illustrated in Figure~\ref{fig:info_trans} requires $10$ edges for full connectivity, yet the butterfly network can accomplish the same with only $8$ edges. The following section details how leveraging the butterfly structure enhances parameter efficiency.

In contrast to the block-diagonal parameterization utilized by OFT to balance between expressive power and regularity, BOFT leverages the butterfly structure to enable a more seamless continuum between the full orthogonal group (which provides high expressivity) and the identity matrix (ensuring regularity). This property allows BOFT to identify a more compact hypothesis space within the set of orthogonal matrices suitable for downstream applications. Given the prominence of butterfly structures in the design of various efficient linear operations, such as the discrete Fourier transform and discrete cosine transform, we contend that our structural approach to parameter efficiency imbues an inherent inductive bias. This inductive bias has the potential to improve generalization and mitigate overfitting. Our rationale stems from the fact that butterfly structures are capable of replicating a wide array of classical linear mappings, a feat not possible with the block-diagonal architecture of OFT. Our key contributions are outlined below:

(\emph{e.g.}, GPT-3 features approximately 175 billion parameters). This exponential scaling renders full-scale training from inception largely inaccessible for most researchers. For continued progress, there is a clear imperative to adapt existing large models for new tasks without starting from scratch. To this end, three predominant paradigms for efficient adaptation have emerged: \emph{model finetuning}~\cite{hulora2022,zaken2022bitfit,guo2020parameter,raffel2020exploring,qiu2023controlling,zhang2022adaptive,chavan2023one}, which selectively updates certain parameters within the established model architecture; \emph{adapter tuning}~\cite{rebuffi2017learning,houlsby2019parameter,pfeiffer2020adapterfusion,lin2020exploring,he2021towards}, which incorporates supplementary trainable modules; and \emph{prompt tuning}~\cite{li2021prefix,lester2021power}, wherein additional prefix tokens are prepended to the input. Among these, model finetuning stands out due to its straightforward implementation and, crucially, because it does not add inference overhead.

Here, $\bm{B}_1(\frac{d}{2})$ and $\bm{B}_2(\frac{d}{2})$ represent butterfly matrices of dimension $\frac{d}{2}$. We introduce the \emph{butterfly component} as $\thickmuskip=2mu \medmuskip=2mu \tilde{\bm{B}}(d,k)=\text{diag}(\bm{B}^F_1(k),\cdots,\bm{B}^F_{d/k}(k))$, a block-diagonal matrix of shape $d\times d$ comprised of blocks of size $k$. Each $\bm{B}^F_i(k)$, as specified in Equation~\ref{eq:but_factor}, refers to a butterfly factor. With these constructs, we can now use the butterfly matrix as a means to parameterize an orthogonal matrix. This construction simply requires that every factor $\tilde{\bm{B}}(d,k)$ in the butterfly matrix $\bm{B}(d)$ be orthogonal for all valid $k$. We first analyze the special case of $\tilde{\bm{B}}(d,2)$, which is a block-diagonal matrix with $2\times 2$ blocks. Orthogonality of $\tilde{\bm{B}}(d,2)$ can be readily ensured by parameterizing each block via the Cayley transform (or equivalently, through $2$-dimensional rotation), consistent with the methods used in \cite{liu2021orthogonal,qiu2023controlling}. Importantly, the arrangement of nonzero entries in each butterfly component corresponds to a permutation of those in $\tilde{\bm{B}}(d,2)$, making it straightforward to ensure that all such components are orthogonally parameterized. In this way, we arrive at an efficient scheme for orthogonal matrix parameterization relying on compositions of $2\times 2$ orthogonal matrices. Extending this idea in the manner of \cite{chen2022pixelated}, we further define a block butterfly component $\tilde{\bm{B}}^b(d,k)$, in which each element of $\bm{d}_i$ is promoted to a $b\times b$ matrix. To ensure that $\tilde{\bm{B}}^b(d,2)$ remains orthogonal, we enforce orthogonality at the level of each $2b\times 2b$ block. For $\tilde{\bm{B}}^b(d,k)$ with $k>2$, the nonzero configuration simply constitutes a blockwise permutation of that in $\tilde{\bm{B}}^b(d,2)$, allowing similar orthogonal parameterizations. By integrating these components, the forward computation in BOFT proceeds as follows:

\textbf{Multiplicative Dropout Mechanism for BOFT.} To regularize the adaptation and curb overfitting, Dropout layers are often employed (as in LoRA~\cite{hulora2022}). However, while conventional Dropout~\cite{srivastava2014dropout} integrates naturally with LoRA's additive updates, it is incompatible with the multiplicative paradigm of BOFT. As a remedy, we introduce a tailored multiplicative Dropout scheme for BOFT: since $\bm{R}(m,b)$ is the product of $m$ butterfly matrices (each permutable into $2b\times2b$ block-diagonal form), the method randomly selects a proportion $p_1$ of butterfly factors and, within each, a proportion $p_2$ of block diagonals to substitute with identity matrices.

\textbf{Butterfly structures}. The radix-2 Cooley-Tukey method recursively breaks down the $N$-point discrete Fourier transform into two smaller $\frac{N}{2}$-point transforms, a process that generates the butterfly structure, expressible as a product of several sparse matrices (collectively referred to as a butterfly matrix). Such matrices have found utility in constructing orthogonal matrix parameterizations for bypassing pivoting during Gaussian elimination and improving computational efficiency~\cite{parker1995random}, as well as for stabilizing recurrent neural network training~\cite{jing2017tunable} and for kernel approximations~\cite{munkhoeva2018quadrature}. The butterfly parameterization has been used to design fast linear transforms~\cite{dao2019learning,dao2019kaleidoscope}. In addition, neural network training efficiency has been boosted through butterfly matrices~\cite{chen2022pixelated,dao2022monarch}. Their applications extend further to areas such as fast matrix-vector multiplication~\cite{michielssen1996multilevel,o2010algorithm}, approximation of data-sparse matrices~\cite{li2015butterfly}, and network broadcasting and transmission~\cite{klasing1994broadcasting,li2003linear,rajasingh2016transmission}. In distinction from prior research, our approach centers on leveraging butterfly structures specifically to advance the parameter efficiency of OFT.

\section{Introduction}

\begin{equation}
    \begin{aligned}\nonumber
    \bm{z}=\big(\bm{R}(m,b)\cdot\bm{W}^0\big)^\top\bm{x},~~~~\text{s.t.}~\bigg{\{}\bm{R}(m,b)=\prod_{i=1}^m \tilde{\bm{B}}^b_{(d,i)}~~\&~~\underbrace{\big(\tilde{\bm{B}}^b_{(d,j)}\big)^\top\tilde{\bm{B}}^b_{(d,j)}=\tilde{\bm{B}}^b_{(d,j)}\big(\tilde{\bm{B}}^b_{(d,j)}\big)^\top\!=\bm{I}_d}_{\forall j\in [1, m]}\bigg{\}},
    \end{aligned}
\end{equation}